\documentclass{article}
\usepackage{graphicx} % Required for inserting images
\usepackage[ruled,vlined,linesnumbered]{algorithm2e}
\usepackage{algpseudocode}
\usepackage{amsthm,amsmath,amssymb}
\usepackage{amssymb,amsfonts}
\usepackage[utf8]{inputenc}
\usepackage{tikz}
\usepackage{natbib,booktabs}
\usepackage{comment}
\usepackage{xcolor}
\usepackage{url}
% \date{1 August, 2024}

\renewcommand{\algorithmicrequire}{\textbf{Input:}}
\renewcommand{\algorithmicensure}{\textbf{Output:}}
\newtheorem{example}{Example}
\newtheorem{theorem}{Theorem}
\newtheorem{definition}{Definition}
\newtheorem{lemma}{Lemma}
\newtheorem{proposition}{Proposition}
\newtheorem{remark}{Remark}

\title{Reserve with score feasibility}
\author{}
\date{\today}

\begin{document}

\maketitle

\section{Model}
\subsection{Basic model (fixed quotas)}
\begin{itemize}
    \item a set of students $I$
    \item There is a single college. 
    The college has the standard category $o$ and 
    a set of reserved categories $T$. Let $V:= T \cup \{o\}$.
\item For each category $c \in V$, $q_c$ is the maximum quota
of $c$. Let $Q=\sum_{c \in V} q_c$
\item For each student $i$ let $V(i)$ denote 
the set of categories that $i$ can apply. 
We assume $o \in V(i)$ holds for each $i\in I$. 
\item a set of contracts 
$X = \{(i,c) \mid i\in I, c \in V(i)\cup\{\emptyset\}\}$.
A contract $(i,c)$ means that student $i$ is accepted to the college under 
category $c$. $(i,\emptyset)$ means $i$ is assigned to an outside option (not assigned to 
any category). 
\item For a set of contracts $Y\subseteq X$, 
Let $Y_i$ denote $\{(i,c) \mid (i,c) \in Y\}$, 
and $Y_c$ denote $\{(i,c) \mid (i,c) \in Y\}$.
\item Each student $i$ has a strict preference $\succ_i$ 
over $X_i$. For brevity, we sometimes write $c \succ_i c'$ 
instead of $(i,c)\succ_i (i,c')$. We write $c \succeq_i c'$ 
when $c \succ_i c'$ or $c=c'$ holds. 
We also use a notation like $Y_i \succ_i c'$, which means
$(i,c) \succ_i (i,c')$ where $Y_i=\{(i,c)\}$.

\item Each student $i$ has her score $\sigma(i) \ge 0$. 
We assume no two students have the same score; we have some tie-breaking rule. 
\end{itemize}

We say $Y\subseteq X$ is a (full) matching 
if $|Y_i|=1$ holds for each student $i \in I$, and
$Y_i \succeq_i \emptyset$ holds. 
We say $Y\subseteq X$ is a partial matching 
if $|Y_i|\le 1$ holds for each student $i \in I$, and
if $Y_i=\{(i,c)\}$, $c \succeq_i \emptyset$ holds. 


\begin{definition}[Feasibility]
A matching $Y$ is feasible if 
    $|Y_c| \leq q_c$ for each $c \in V$.
\end{definition}
For a matching $Y$, 
let $R_o(Y)$ denote $\{i\in I \mid (i,\emptyset) \in Y, o\text{ is acceptable for }i\}$, 
i.e., $R_o(Y)$ is the set of students who are assigned to the outside option although they prefer
$o$ over the outside option. 
Then, $e_o(Y)$ is defined as follows.
$e_o(Y)=0$ if $|R_o(Y)|=0$. Otherwise, 
$e_o(Y) = \max_{i \in R_o(Y)}\sigma(i)$.
Intuitively, $e_o(Y)$ is the score of the top student who is rejected from the college 
although she hopes to be assigned to the standard category. 

For each category $t\in T$, its 
acceptable margin $\delta_t \geq 0$ is given. 
\begin{definition}[Score-feasibility]
We say $Y$ is score-feasible if for each $i$ such that 
$(i,t) \in Y, t \in T$, $\sigma(i) \ge e_o(Y) - \delta_t$ holds. 
\end{definition}
In words, score-feasibility guarantees the difference between 
(i) the 
highest score of the student who is rejected from the college (although 
$o$ is acceptable for her),
and (ii) the lowest score of the student who is accepted in 
each category $t$, is at most $\delta_t$. 

An alternative way to define score-feasibility is using the 
\emph{lowest score of the accepted student} in category $o$, 
rather than the \emph{highest score of the rejected student}.
These two definitions look similar, and their difference becomes 
small in practice when many borderline students exist. 
We use the current definition since the following property holds. 

\begin{proposition}
    Assume matching $Y$ is feasible and score-feasible. 
    Choose $i$ as $\arg\min_{(i,c)\in Y} \sigma(i)$.
    Then, $Y \setminus \{(i,c)\}$ is also feasible and score-feasible. 
\end{proposition}
In words, if $Y$ is feasible and score-feasible, 
the matching obtained by removing the lowest score student from $Y$ is also 
feasible and score-feasible. This property is crucial in our proposed mechanisms. 

This property does not hold if we use the lowest score of the accepted student.
Assume $T=\{t\}$, three students $i_1$, where $V(i_1)=\{o\}$, $\sigma(i_1)=100$, 
$i_2$, where $V(i_2)=\{o,t\}$, $\sigma(i_2)=80$, 
and $i_3$, where $V(i_3)=\{o\}$, $\sigma(i_3)=60$. Let $\delta_t = 10$. 
Assume $Y=\{(i_1, o), (i_2,t), (i_3, o)\}$. 
This is feasible and score-feasible. 
If we remove the lowest-score student, i.e., $i_3$, 
we obtain $Y'=\{(i_1, o), (i_2,t)\}$. 
This is feasible and score-feasible in our definition; 
all accepted students have higher scores than the rejected student $i_3$. 
If we require that the score difference between the lowest-score accepted student 
in $o$ and the lowest-score accepted student in $t$ is at most $\delta_t$, 
$Y$ is feasible and score-feasible, but $Y'$ is not score-feasible, since 
the lowest score of the accepted student in $o$, i.e., $i_1$, is 100, 
while the score of $i_2$ is 80.

\begin{definition}[fairness]
    Student $i$ has justified envy toward $j$ in matching $Y$
if $Y_j=\{(j, c)\}, c \in V(i)$, 
$c \succ_i Y_i$, and $\sigma(i)>\sigma(j)$ hold.
Matching is fair if no student has justified envy.
\end{definition}

\begin{definition}[Non-wastefulness]
In matching $Y$, student $i$ claims an empty seat of category $c$,
if $c \in V(i)$, $i$ prefers $c$ over $Y_i$, 
and moving $i$ to $c$ is feasible and score-feasible. 
A matching is non-wasteful if no student claims an empty seat.
\end{definition}

\begin{definition}[Pareto efficiency]
We say a feasible and score-feasible matching is Pareto efficient (PE), 
if it is not Pareto dominated by another feasible and score-feasible matching.
\end{definition}
PE implies non-wastefulness. 

\begin{definition}[SSD]
Score-based Serial dictatorship mechanism (SSD) is defined 
as follows:
\begin{enumerate}
    \item 
Choose student $i$ one by one in decreasing order of score.
\item Assign $i$ to her most favorite category (including the outside option $\emptyset$) 
s.t. the resulting (partial)
   matching is feasible and score-feasible. More specifically, 
   we set $e_o(Y)=0$ until the first student $i$ is assigned 
to the outside option while $i$ prefers $o$ over the outside option. 
After that, $e_o(Y)=\sigma(i)$. 
\end{enumerate}   
\end{definition}

\begin{theorem}
SSD is feasible, score-feasible, fair, PE, and strategy-proof.     
\end{theorem}
\begin{proof}[Proof sketch:]
Feasibility and score-feasibility are clear from construction. 
Since it is a serial dictatorship mechanism, PE and strategy-proofness are also clear. 
Also, since the selection order of students is based on the score, 
fairness is also clear. 
\end{proof}


\subsection{Flexible quota model}
Feasibility is modified as follows. 
\begin{definition}[Feasibility]
A matching $Y$ is feasible if 
    $|Y_c| \leq q_c$ for each $c \in T$, 
    and $\sum_{c \in V} |Y_c| \leq Q$ hold. 
\end{definition}

The definition of $e_o(Y)$ and score-feasibility are the same 
as the fixed quota model. 

\begin{definition}[Reserved category protection]
For each student $i$, let $\text{top}_i$ denote her top category in $V(i) \cup \{\emptyset\}$.
For each category $c\in T$, let $d_c$ denote $\min{\{q_c, |\{i \mid \text{top}_i =c\}|\}}$.
We say matching $Y$ satisfies reserved category protection 
if $|Y_o|\leq q_o + \sum_{c \in T} (q_c - \max\{|Y_c|, d_c\})$ holds. 
\end{definition}

\begin{theorem}
There is a case where no matching 
simultaneously satisfies feasibility, score-feasibility, non-wastefulness, and
reserved category protection.
\end{theorem}
\begin{proof}
Consider the case where $T=\{t\}$, $q_o=q_t=1$. 
There are three students: $i_1$, where $\sigma(i_1)=100$, 
$V(i_1)=\{o\}$, and $o\succ_{i_1} \emptyset$.
$i_2$, where $\sigma(i_2)=90$, 
$V(i_2)=\{o\}$, and $o\succ_{i_2} \emptyset$.
$i_3$, where $\sigma(i_3)=60$, 
$V(i_3)=\{o, t\}$, and 
$t \succ_{i_3} o \succ_{i_3} \emptyset$. 
$\delta_t=10$. 

By reserved category protection, one seat must be reserved for category $t$. 
Thus, at least one of $i_1$ or $i_2$ is assigned to the outside option. 
In either case, 
assigning $i_3$ to $t$ violates score-feasibility. 
Thus, no student is assigned to $t$ and at most one student can be assigned to $o$. 
However, either assigning no student to $o$, or assigning one student to $o$, violates
non-wastefulness (since no student is assigned to the reserved category, moving one unassigned 
student to $o$ will satisfy feasibility and score-feasibility). 
\end{proof}

Since non-wastefulness is too strong, let us introduce a weaker efficiency requirement.
    \begin{definition}[Weak non-wastefulness]
In matching $Y$, student $i$ strongly claims an empty seat of category $c$,
if $c \in V(i)$, $i$ prefers $c$ over $Y_i$, 
and moving $i$ to $c$ is feasible, score-feasible, and satisfies reserved category protection.
A matching is weakly non-wasteful if no student strongly claims an empty seat.
\end{definition}

However, we still obtain an impossibility result if we additionally require fairness and 
strategy-proofness. 

\begin{theorem}
No mechanism 
simultaneously satisfies feasibility, score-feasibility, fairness, weak non-wastefulness, 
reserved category protection, and strategy-proofness in general. 
\end{theorem}
\begin{proof}
Consider the case where $T=\{t\}$, $q_o=q_t=1$. 
There are three students: $i_1$, where $\sigma(i_1)=100$, 
$V(i_1)=\{o\}$, and $o\succ_{i_1} \emptyset$.
$i_2$, where $\sigma(i_2)=90$, 
$V(i_2)=\{o\}$, and $o\succ_{i_2} \emptyset$.
$i_3$, where $\sigma(i_3)=80$, 
$V(i_3)=\{o, t\}$, and 
$o \succ_{i_3} t \succ_{i_3} \emptyset$. 
$\delta_t=10$. 
Assume all students declare their true preferences. 
Then, there is no need to protect reserved categories. By fairness and weak non-wastefulness, 
$i_1$ and $i_2$ are assigned to $o$, while $i_3$ is assigned to the outside option. 

Next, assume $i_3$ misreports her preference as 
$t \succ'_{i_3} o \succ'_{i_3} \emptyset$. 
Then, one seat must be protected to $t$. 
The only assignment that satisfies the required conditions is assigning $i_1$ to $o$ and 
$i_3$ to $t$. Thus, $i_3$ obtains a strictly better outcome by misreporting; strategy-proofness
is violated. 
\end{proof}

\begin{definition}[Reserved category protecting SSD]
Reserved category protecting SSD is defined as follows. 
Run SSD, requiring feasibility, score-feasibility, and reserved category protection based on 
$d_c$. More specifically, the number of students assigned to $o$ is at most
$Q-\sum_{c\in T}\max{\{|Y_c|, d_c\}}$. Also, we set $e_o(Y)=0$ until the first student $i$ is assigned 
to the outside option while $i$ prefers $o$ over the outside option. 
After that, $e_o(Y)=\sigma(i)$. 
\end{definition}

\begin{theorem}
Reserved category protecting SSD
satisfies feasibility, score-feasibility, fairness, weak non-wastefulness, 
and reserved category protection. Also, it is PE within all matchings satisfying 
feasibility, score-feasibility, and reserved category protection. 
\end{theorem}
Note that Reserved category protecting SSD is not strategy-proof, 
by the above impossibility theorem.

\begin{definition}[flexible SSD]
Flexible SSD is defined as follows. 
Run SSD, requiring feasibility based on flexible quotas.
More specifically, for the current partial matching $Y$, if 
total number of students assigned is equal to $Q$
(except for who are assigned to the outside option), a student can be assigned only to 
the outside option. 
Otherwise, she can be assigned to $o$. Also, if $|Y_c|<q_c$ for $c\in T$ and $c\in V(i)$, she can 
be assigned to $c$.
\end{definition}

\begin{theorem}
Flexible SSD
satisfies feasibility, score-feasibility, fairness, PE, and strategy-proofness.
\end{theorem}
Note that flexible SSD does not satisfy reserved category protection in general.

\begin{definition}[fixed quota SSD]
Fixed quota SSD is defined as follows. 
Set $q_o = Q - \sum_{c \in T} q_c$. 
Run SSD with $q_o$ and $q_c$ for each $c \in T$.
\end{definition}
\begin{theorem}
Fixed quota SSD satisfies feasibility, score-feasibility, fairness, reserved category protection, and strategy-proofness.      
\end{theorem}
Note that fixed quota SSD does not satisfy weak non-wastefulness, by the above impossibility theorem. 

% shimamura added on 10/5 about eligibility based ssd

\subsection{Eligibility-based reserved category protection}
Reserved category protecting SSD is not strategy-proof.
We think the reason is that the protection level $d_c$ depends on the reports of students.
We show two examples. Below, we write $d^{\mathrm{top}}_c$ for $d_c$.

\begin{example}[Misreporting the top choice]
In the proof of Theorem 3, $d^{\mathrm{top}}_t$ changes from $0$ to $1$ when $i_3$
reports $t$ as her top choice.
Then $i_3$ is assigned to $t$ instead of the outside option.
\end{example}

\begin{example}[Hiding an eligible category]
Assume the protection level is decided by the categories that each student declares.
More specifically, student $i$ declares $\tilde{V}(i)\subseteq V(i)$, and we use
$\tilde{d}_c=\min\{q_c, |\{i \mid c \in \tilde{V}(i)\}|\}$ in
Reserved category protecting SSD.
Consider the case where $T=\{t\}$, $q_o=q_t=1$.
There are two students: $i_1$, where $\sigma(i_1)=100$, $V(i_1)=\{o\}$, and
$o\succ_{i_1}\emptyset$.
$i_2$, where $\sigma(i_2)=90$, $V(i_2)=\{o,t\}$, and
$o\succ_{i_2} t\succ_{i_2}\emptyset$.

Assume $i_2$ declares $\{o,t\}$. Then $\tilde{d}_t=1$, so only one student can be assigned to $o$.
$i_1$ is assigned to $o$ and $i_2$ is assigned to $t$.

Next, assume $i_2$ declares $\{o\}$. Then $\tilde{d}_t=0$, so two students can be assigned to $o$.
$i_2$ is assigned to $o$.
Thus, $i_2$ obtains a strictly better outcome by hiding $t$.
\end{example}

In both examples, the protection level changes by the report of a student.
So we decide the protection level by using only $V(i)$.
We assume $V(i)$ is checked by the college, and student $i$ cannot change it.
Each student reports only her preference.

\begin{definition}[Eligibility-based reserved category protection]
For each category $c\in T$, let $d^V_c$ denote $\min\{q_c, |\{i \mid c\in V(i)\}|\}$.
We say matching $Y$ satisfies eligibility-based reserved category protection
if $|Y_o|\leq q_o+\sum_{c\in T}(q_c-\max\{|Y_c|, d^V_c\})$ holds.
\end{definition}

\begin{proposition}
If matching $Y$ satisfies eligibility-based reserved category protection,
then $Y$ satisfies reserved category protection.
\end{proposition}
\begin{proof}
If $\text{top}_i=c$ and $c\in T$, then $c\in V(i)$ holds.
Thus, $d^{\mathrm{top}}_c\le d^V_c$ holds for each $c\in T$.
Then, $\max\{|Y_c|, d^{\mathrm{top}}_c\}\le\max\{|Y_c|, d^V_c\}$ holds, so
\[
|Y_o|\le q_o+\sum_{c\in T}(q_c-\max\{|Y_c|, d^V_c\})
\le q_o+\sum_{c\in T}(q_c-\max\{|Y_c|, d^{\mathrm{top}}_c\}).
\]
\end{proof}

Next, we show that this model can be reduced to the fixed quota model.
Let $q'_o := q_o+\sum_{c\in T}(q_c-d^V_c)$.
$q'_o$ is a constant that is decided only by the quotas and $V(i)$.

\begin{lemma}\label{lem:reduction}
Matching $Y$ satisfies feasibility (of the flexible quota model) and
eligibility-based reserved category protection if and only if
$|Y_o|\le q'_o$ and $|Y_c|\le q_c$ for each $c\in T$ hold.
\end{lemma}
\begin{proof}
In both directions, $|Y_c|\le q_c$ holds for each $c\in T$.
Also, if $(i,c)\in Y$, then $c\in V(i)$. So $|Y_c|\le |\{i \mid c\in V(i)\}|$ holds.
Thus, $|Y_c|\le d^V_c$, i.e., $\max\{|Y_c|, d^V_c\}=d^V_c$.
Then, eligibility-based reserved category protection is the same as $|Y_o|\le q'_o$.

It remains to show $\sum_{c\in V}|Y_c|\le Q$ when $|Y_o|\le q'_o$ holds.
This holds since
$\sum_{c\in V}|Y_c|\le q'_o+\sum_{c\in T}d^V_c=q_o+\sum_{c\in T}q_c=Q$.
\end{proof}

We define the efficiency requirements based on $d^V_c$ in the same way.
\begin{definition}[Weak non-wastefulness ($d^V$)]
In matching $Y$, student $i$ strongly claims ($d^V$) an empty seat of category $c$,
if $c\in V(i)$, $i$ prefers $c$ over $Y_i$, and moving $i$ to $c$ is feasible, score-feasible,
and satisfies eligibility-based reserved category protection.
A matching is weakly non-wasteful ($d^V$) if no student strongly claims ($d^V$) an empty seat.
\end{definition}

\begin{definition}[Constrained PE ($d^V$)]
We say a matching that satisfies feasibility, score-feasibility, and eligibility-based
reserved category protection is constrained PE ($d^V$),
if it is not Pareto dominated by another matching that satisfies these three conditions.
\end{definition}
To distinguish, we call the requirements in Theorem 4 weak non-wastefulness
($d^{\mathrm{top}}$) and constrained PE ($d^{\mathrm{top}}$).

\begin{definition}[Eligibility-based SSD]
Eligibility-based SSD is defined as follows.
Run SSD of the fixed quota model, where the quota of $o$ is $q'_o$ and
the quota of each $c\in T$ is $q_c$.
\end{definition}

\begin{theorem}
Eligibility-based SSD satisfies feasibility, score-feasibility, fairness,
eligibility-based reserved category protection, weak non-wastefulness ($d^V$),
and strategy-proofness. Also, it is constrained PE ($d^V$).
\end{theorem}
\begin{proof}
Let $Y$ be the obtained matching.
Eligibility-based SSD is SSD of the fixed quota model with $q'_o$, so we can use Theorem 1.

\textbf{Feasibility and eligibility-based reserved category protection.}
By Theorem 1, $|Y_o|\le q'_o$ and $|Y_c|\le q_c$ for each $c\in T$ hold.
Thus, by Lemma~\ref{lem:reduction}, $Y$ satisfies both.

\textbf{Score-feasibility and fairness.}
The definitions are the same as the fixed quota model, so they hold by Theorem 1.

\textbf{Constrained PE ($d^V$).}
Assume matching $Z$ satisfies feasibility, score-feasibility, and eligibility-based
reserved category protection, and $Z$ Pareto dominates $Y$.
By Lemma~\ref{lem:reduction}, $Z$ is feasible in the fixed quota model with $q'_o$.
This contradicts the fact that SSD is PE (Theorem 1).

\textbf{Weak non-wastefulness ($d^V$).}
Assume student $i$ strongly claims ($d^V$) an empty seat of $c$.
Let $Z$ be the matching obtained by moving $i$ to $c$.
Then $Z$ satisfies the three conditions, $i$ prefers $Z$ over $Y$, and the other students
are the same. Thus, $Z$ Pareto dominates $Y$.
This contradicts constrained PE ($d^V$).

\textbf{Strategy-proofness.}
$q'_o$ is decided only by the quotas and $V(i)$, so it does not change by the reports of students.
Thus, strategy-proofness holds by Theorem 1.
\end{proof}

\begin{example}
Consider the case in the proof of Theorem 3. Since $d^V_t=1$, we have $q'_o=1$.
First, $i_1$ is assigned to $o$.
Next, $i_2$ is assigned to the outside option since $o$ is full, and $e_o(Y)=90$.
Finally, $i_3$ is assigned to $t$ since $o$ is full and $80\ge 90-\delta$ holds.
The outcome is the same when $i_3$ reports $t\succ'_{i_3}o\succ'_{i_3}\emptyset$.
So $i_3$ cannot obtain a better outcome by misreporting.

However, in this outcome, $i_3$ prefers $o$ over $t$.
Since $d^{\mathrm{top}}_t=0$, moving $i_3$ to $o$ satisfies feasibility,
score-feasibility, and reserved category protection.
Thus, Eligibility-based SSD is not weakly non-wasteful ($d^{\mathrm{top}}$).
It is also not non-wasteful, not PE, and not constrained PE ($d^{\mathrm{top}}$).
\end{example}

The properties of the three mechanisms are as follows.
``---'' means that the mechanism does not satisfy the corresponding protection.
\begin{center}
\begin{tabular}{lccc}
\toprule
 & RCP SSD & Flexible SSD & Eligibility-based SSD\\
\midrule
Feasibility & \checkmark & \checkmark & \checkmark\\
Score-feasibility & \checkmark & \checkmark & \checkmark\\
Fairness & \checkmark & \checkmark & \checkmark\\
Non-wastefulness & $\times$ & \checkmark & $\times$\\
Weak non-wastefulness ($d^{\mathrm{top}}$) & \checkmark & \checkmark & $\times$\\
Weak non-wastefulness ($d^V$) & \checkmark & \checkmark & \checkmark\\
PE & $\times$ & \checkmark & $\times$\\
Constrained PE ($d^{\mathrm{top}}$) & \checkmark & --- & $\times$\\
Constrained PE ($d^V$) & --- & --- & \checkmark\\
Reserved category protection & \checkmark ($d^{\mathrm{top}}$) & $\times$ & \checkmark ($d^V$)\\
Strategy-proofness & $\times$ & \checkmark & \checkmark\\
\bottomrule
\end{tabular}
\end{center}

\subsection{Proofs}
\begin{proof}[Proof of Theorem 1]
Let $Y$ be the matching obtained by SSD.
For student $i$, let $Y^i$ denote the partial matching just before $i$ is chosen.

First, we check $e_o(Y)$.
Students are chosen in decreasing order of score.
So the first student who is assigned to the outside option although $o$ is acceptable
has the highest score in $R_o(Y)$.
Thus, the value set in SSD is equal to $e_o(Y)$, and it does not change after it is set.

\textbf{Feasibility.} Clear from the construction.

\textbf{Score-feasibility.}
Assume $(i,c)\in Y$ and $c\in T$.
If $e_o(Y)$ is already set when $i$ is chosen, $\sigma(i)\ge e_o(Y)-\delta$ holds
from the construction.
Otherwise, all students in $R_o(Y)$ are chosen after $i$, so $e_o(Y)<\sigma(i)$ holds.

\textbf{Fairness.}
Assume $i$ has justified envy toward $j$, where $Y_j=\{(j,c)\}$.
Since $\sigma(i)>\sigma(j)$, $i$ is chosen before $j$.
Since $i$ prefers $c$ over $Y_i$, $i$ could not choose $c$. There are two cases.
\begin{enumerate}
\item $c$ is full when $i$ is chosen.
Then $c$ is also full when $j$ is chosen, so $j$ cannot be assigned to $c$.
\item $c\in T$ and $\sigma(i)<e_o(Y)-\delta$.
Then $\sigma(j)<\sigma(i)<e_o(Y)-\delta$, so $j$ cannot be assigned to $c$.
\end{enumerate}
Both cases contradict $Y_j=\{(j,c)\}$.

\textbf{PE.}
Assume matching $Z$ is feasible and score-feasible, and $Z$ Pareto dominates $Y$.
Let $i$ be the student with the highest score among students s.t. $Z_i\neq Y_i$.
Then $i$ prefers $Z_i$ over $Y_i$. Let $Z_i=\{(i,c)\}$.
$c\neq\emptyset$ holds, since $i$ can always choose the outside option in SSD.
The students chosen before $i$ have the same assignment in $Y$ and $Z$,
so $Y^i\cup\{(i,c)\}\subseteq Z$ holds.
Since $Z$ is feasible, $c$ is not full when $i$ is chosen.
Also, assume $c\in T$.
The students in $R_o(Y)$ chosen before $i$ are also in $R_o(Z)$, so the value of $e_o$
when $i$ is chosen is at most $e_o(Z)$.
Since $Z$ is score-feasible, $\sigma(i)\ge e_o(Z)-\delta$ holds.
Therefore, $i$ could choose $c$ in SSD.
This contradicts the fact that $i$ is assigned to her most favorite category.

\textbf{Strategy-proofness.}
The categories that $i$ can choose are decided by $Y^i$ and the value of $e_o$ at that time.
They depend only on the reports of students whose scores are higher than $\sigma(i)$.
Thus, $i$ cannot change them by her report, and reporting truthfully is the best for $i$.
\end{proof}

\begin{proof}[Proof of Theorem 5]
The proof is almost the same as Theorem 1.
The only difference is the condition that $i$ can choose $c$:
the total number of assigned students is less than $Q$ (and $|Y_c|<q_c$ if $c\in T$).
In the proof of fairness, if $i$ could not choose $c$ by this condition, $j$ cannot choose
$c$ either, since the number of assigned students does not decrease.
In the proof of PE, $Y^i\cup\{(i,c)\}\subseteq Z$ and $Z$ is feasible,
so $i$ could choose $c$.
The other parts are the same.
\end{proof}

\begin{proof}[Proof of Theorem 4]
Since $Q=q_o+\sum_{c\in T}q_c$, reserved category protection is equivalent to
\begin{equation}\label{eq:rcp}
|Y_o|+\sum_{c\in T}\max\{|Y_c|, d_c\}\le Q.
\end{equation}
The left-hand side does not decrease when a student is assigned to a category.

\textbf{Feasibility, score-feasibility, and reserved category protection.}
The empty matching satisfies (\ref{eq:rcp}) since $\sum_{c\in T}d_c\le Q$.
The mechanism keeps (\ref{eq:rcp}) and $|Y_c|\le q_c$ for $c\in T$, and
(\ref{eq:rcp}) implies $\sum_{c\in V}|Y_c|\le Q$.
Score-feasibility is the same as Theorem 1.

\textbf{Fairness.}
The proof is the same as Theorem 1, except for the following case:
$i$ could not choose $c$ because (\ref{eq:rcp}) is violated.
In this case, since the left-hand side of (\ref{eq:rcp}) does not decrease,
$j$ cannot choose $c$ either.

\textbf{PE within matchings satisfying the three conditions.}
Assume matching $Z$ satisfies feasibility, score-feasibility, and reserved category protection,
and $Z$ Pareto dominates $Y$.
We choose $i$ and $c$ in the same way as Theorem 1.
Since $Y^i\cup\{(i,c)\}\subseteq Z$ and $Z$ satisfies (\ref{eq:rcp}),
$Y^i\cup\{(i,c)\}$ also satisfies (\ref{eq:rcp}).
The other conditions are the same as Theorem 1, so $i$ could choose $c$, a contradiction.

\textbf{Weak non-wastefulness.}
Assume student $i$ strongly claims an empty seat of $c$.
The matching obtained by moving $i$ to $c$ satisfies the three conditions and
Pareto dominates $Y$. This contradicts the above.
\end{proof}

\end{document}

